\section{Discussion}
\subsection{Result Analysis}
The most important observation is that the nominal source$\rightarrow$target transfer is easier than expected from mass alone. A source-trained SAC policy reaches 96\% success on the 5 kg target object, only slightly below its 98\% source success. This suggests that the current \texttt{PandaPush-v3} setup is not very sensitive to the fixed 1 kg to 5 kg mass shift. The task has a fixed goal, a short object-goal distance distribution, dense rewards during training, and a powerful continuous-control algorithm. These factors can hide part of the intended reality gap.

The harder randomized evaluation is therefore more informative. When target mass is sampled in $[5,10]$ kg and friction is randomized, the non-randomized SAC policy falls to 76\% success. This indicates that friction and heavier masses expose failures that are not visible in the nominal target evaluation. The UDR policy trained with mass and friction randomization performs worse than the nominal SAC policy on the easy fixed evaluation, but better on the hard evaluation. This is the expected robustness trade-off: UDR may sacrifice performance on the nominal dynamics in exchange for a policy that is less specialized.

SAC appears more suitable than PPO for this repository state. The available logs show high SAC training success after 500k steps, while PPO runs are less consistent and require careful treatment of vector normalization. This is coherent with the algorithms: SAC is off-policy and can reuse data through its replay buffer, whereas PPO is on-policy and discards rollouts after each update. For goal-conditioned pushing, replay-based learning is especially useful because informative transitions can be reused many times.

\subsection{Reward Shaping}
The reward-shaping wrapper affects training but not the evaluation success metric. The time penalty encourages short solutions, and the close-goal bonus increases the learning signal once the object reaches the success threshold. For SAC, this can increase success rate indirectly by making successful behavior easier to learn early in training. However, it also makes training return less comparable to raw evaluation return. This is why the evaluation script reports both return and success rate on the unwrapped environment.

\subsection{Methodological Critique}
The main limitation is that the current result set is not yet statistically complete. Most numbers are single-seed evaluations with 50 episodes. They are useful for debugging and for identifying trends, but the final report should include mean and standard deviation across seeds where possible.

The domain-randomization implementation is also intentionally simple. UDR uses hand-selected uniform ranges, so its behavior depends strongly on those ranges. If the training range matches the evaluation range too closely, the comparison becomes too generous. The current setup avoids the most direct form of this issue by using conservative training friction ranges and a harder evaluation range, but a final study should report the exact train and test distributions clearly.

ADR is currently mass-only and uses a simple success-rate heuristic. It does not adapt friction, center of mass, action noise, observation noise, or target sampling. Therefore, any ADR result should be interpreted as a curriculum over object mass rather than as a full automatic domain-randomization method. In addition, the source-domain success rate may be a weak curriculum signal for transfer: a policy can be successful in the randomized training domain while still failing on the target distribution.

Finally, PPO evaluation needs care because some saved PPO models were trained behind \texttt{VecNormalize}. A valid deterministic PPO evaluation must load the corresponding normalization statistics; otherwise the policy sees observations on a scale different from training. This is marked as TODO instead of reporting unreliable numbers.
